% =====================================================================
% sec-appendix.tex — 附录
% 附录 A 数据集与样例
% 附录 B 核心算法实现
% 附录 C 15 分箱与样本量收敛数值明细表
% 五号字体，不编号节，写入目录
% =====================================================================
\phantomsection
\addcontentsline{toc}{section}{附录}
\section*{附录}

\phantomsection
\addcontentsline{toc}{subsection}{附录 A 数据集与预测样本样例}
\subsection*{附录 A 数据集与预测样本样例}

\textbf{A.1 数据集与基准模型著录。}
实证研究所用 CIFAR-10 数据集源自多伦多大学官方图像基准（Krizhevsky \& Hinton, 2009~\cite{krizhevsky2009learning}），测试集包含 10 个类别均衡的 10,000 张 $32 \times 32$ 彩色图像。预训练残差网络主干（ResNet-20 与 ResNet-56，He 等, 2016~\cite{he2016deep}）权重取自 PyTorch 官方基准模型仓。在分层抽样协议下，校准拟合集（$n_{\mathrm{cal}}=2000$）与独立评估集（$n_{\mathrm{eval}}=8000$）严格互斥隔离。

\textbf{A.2 独立评估集前 10 个测试样本预测样例。}
表~\ref{tab:app_sample_data}列出了独立评估集前 10 个样本在 ResNet-20 下校准前后的预测类别与最大置信度对照。从中可见，单参数 Temperature Scaling 在保持预测类别严格不变的前提下，使高置信度平滑合理下调。

\begin{table}[htbp]
\centering\small
\caption{独立评估集前 10 个样本在 ResNet-20 模型下校准前后预测结果对照}\label{tab:app_sample_data}
\resizebox{\textwidth}{!}{%
\begin{tabular}{cccccc}
\toprule
样本序号 & 真实类别 $Y_i$ & 未校准预测 $\hat Y_i$ & 未校准置信度 $\hat P_i$ & TS 校准预测 $\hat Y_i^{\mathrm{ts}}$ & TS 校准置信度 $\hat P_i^{\mathrm{ts}}$ \\
\midrule
1 & 9 (卡车) & 9 (卡车) & 1.0000 & 9 (卡车) & 0.9968 \\
2 & 5 (狗) & 5 (狗) & 1.0000 & 5 (狗) & 0.9996 \\
3 & 7 (马) & 7 (马) & 1.0000 & 7 (马) & 0.9996 \\
4 & 0 (飞机) & 0 (飞机) & 1.0000 & 0 (飞机) & 0.9991 \\
5 & 1 (汽车) & 1 (汽车) & 1.0000 & 1 (汽车) & 0.9995 \\
6 & 2 (鸟) & 2 (鸟) & 1.0000 & 2 (鸟) & 0.9998 \\
7 & 5 (狗) & 5 (狗) & 0.9963 & 5 (狗) & 0.9571 \\
8 & 8 (船) & 8 (船) & 1.0000 & 8 (船) & 0.9982 \\
9 & 6 (青蛙) & 2 (鸟) & 0.7927 & 2 (鸟) & 0.6872 \\
10 & 7 (马) & 7 (马) & 1.0000 & 7 (马) & 0.9979 \\
\bottomrule
\end{tabular}%
}
\end{table}


\phantomsection
\addcontentsline{toc}{subsection}{附录 B 核心算法极简实现}
\subsection*{附录 B 核心算法极简实现}

本附录节选实验程序（\texttt{run\_calibration\_experiment.py}）中的核心数值实现，包括用于求解逆温度极大似然估计的阻尼 Newton--Raphson 迭代、观测 Fisher 信息量计算以及针对原假设 $H_0:\beta=1$ 的 Wald、Score 与 LRT 检验函数。

\begin{center}\footnotesize
\begin{verbatim}
import numpy as np

def fit_temperature_mle(logits, labels, max_iter=50, tol=1e-12):
    """极大似然阻尼牛顿法求解逆温度 beta 与 Fisher 信息量 (核心极简实现)"""
    n, c = logits.shape; beta = 1.0; y_onehot = np.eye(c)[labels]
    for _ in range(max_iter):
        shift = np.max(beta * logits, axis=1, keepdims=True)
        p = np.exp(beta * logits - shift)
        p /= np.sum(p, axis=1, keepdims=True)
        exp_z = np.sum(p * logits, axis=1)                          # E[Z_i]
        score = np.sum(logits[np.arange(n), labels] - exp_z)        # U(beta)
        var_z = np.sum(p * (logits**2), axis=1) - (exp_z ** 2)
        fisher = np.sum(var_z)                                      # I_n(beta)
        step = score / max(fisher, 1e-12); gamma = 1.0
        while beta + gamma * step < 0.05 and gamma > 1e-4: gamma *= 0.5
        beta_new = max(beta + gamma * step, 0.05)
        if abs(beta_new - beta) < tol: break
        beta = beta_new
    T_hat = 1.0 / beta; se_T = (T_hat ** 2) / np.sqrt(fisher)      # Delta SE
    return beta, T_hat, se_T, fisher

def compute_hypothesis_tests(logits, labels, beta_hat, fisher_hat):
    """计算针对 H0: beta=1 的三大检验统计量 (Wald, Score, LRT)"""
    n, c = logits.shape; wald = ((beta_hat - 1.0) ** 2) * fisher_hat
    p1 = np.exp(logits - np.max(logits, axis=1, keepdims=True))
    p1 /= np.sum(p1, axis=1, keepdims=True)
    exp_z1 = np.sum(p1 * logits, axis=1)
    u1 = np.sum(logits[np.arange(n), labels] - exp_z1)
    i1 = np.sum(np.sum(p1 * (logits**2), axis=1) - (exp_z1 ** 2))
    score = (u1 ** 2) / max(i1, 1e-12)
    def nll(b):
        sh = np.max(b * logits, axis=1, keepdims=True)
        p = np.exp(b * logits - sh)
        return np.sum(np.log(np.sum(p, axis=1)) + sh.squeeze()
                      - b * logits[np.arange(n), labels])
    lrt = 2.0 * (nll(1.0) - nll(beta_hat))
    return {"Wald": wald, "Score": score, "LRT": lrt}
\end{verbatim}
\end{center}

\phantomsection
\addcontentsline{toc}{subsection}{附录 C 评估集 15 分箱与样本量收敛数值明细表}
\subsection*{附录 C 评估集 15 分箱与样本量收敛数值明细表}

\textbf{C.1 评估集 15 分箱校准数值明细表。}
表~\ref{tab:app_bins}汇报了独立评估集（$n_{\mathrm{eval}}=8000$）在 Temperature Scaling 校准后各分箱的样本容量、经验准确率及 ResNet-56 的平均置信度（对应第~\ref{sec:empirical} 节可靠性图）。

\begin{table}[htbp]
\centering\small
\caption{独立评估集（$n_{\mathrm{eval}}=8000$）Temperature Scaling 校准后 15 分箱详细统计表}\label{tab:app_bins}
\resizebox{\textwidth}{!}{%
\begin{tabular}{ccccccc}
\toprule
分箱编号 & 置信度区间 & 样本量 (R20) & 准确率 (R20) & 样本量 (R56) & 准确率 (R56) & 平均置信度 (R56) \\
\midrule
1--3 & $[0.00, 0.20]$ & 0 & --- & 0 & --- & --- \\
4 & $(0.20, 0.27]$ & 1 & $0.0000$ & 1 & $0.0000$ & $0.2649$ \\
5 & $(0.27, 0.33]$ & 27 & $0.4815$ & 12 & $0.2500$ & $0.2930$ \\
6 & $(0.33, 0.40]$ & 52 & $0.3654$ & 27 & $0.3333$ & $0.3712$ \\
7 & $(0.40, 0.47]$ & 91 & $0.4176$ & 36 & $0.5000$ & $0.4457$ \\
8 & $(0.47, 0.53]$ & 149 & $0.6040$ & 81 & $0.4321$ & $0.5047$ \\
9 & $(0.53, 0.60]$ & 153 & $0.5817$ & 93 & $0.5591$ & $0.5635$ \\
10 & $(0.60, 0.67]$ & 146 & $0.6849$ & 88 & $0.5682$ & $0.6311$ \\
11 & $(0.67, 0.73]$ & 208 & $0.6779$ & 111 & $0.6126$ & $0.6977$ \\
12 & $(0.73, 0.80]$ & 188 & $0.6862$ & 114 & $0.6842$ & $0.7682$ \\
13 & $(0.80, 0.87]$ & 267 & $0.7940$ & 201 & $0.7463$ & $0.8371$ \\
14 & $(0.87, 0.93]$ & 449 & $0.8597$ & 351 & $0.8348$ & $0.9057$ \\
15 & $(0.93, 1.00]$ & 6269 & $0.9880$ & 6885 & $0.9856$ & $0.9882$ \\
\bottomrule
\end{tabular}%
}
\end{table}


\textbf{C.2 样本量收敛轨迹数值明细表。}
表~\ref{tab:app_convergence}汇报了校准集样本量从 200 扩大至 4000 时，50 次重复采样的温度估计均值 $\bar T$、经验标准差（Empirical SD）及理论 Delta 渐近标准误（Delta SE），完整反映了 $\sqrt{n}$ 相合速率与有限总体修正机制。

\begin{table}[htbp]
\centering\small
\caption{校准集样本量收敛轨迹明细表（基于 50 次独立重复分层抽样）}\label{tab:app_convergence}
\resizebox{\textwidth}{!}{%
\begin{tabular}{ccccccc}
\toprule
\multirow{2}{*}{校准样本量 $n_{\mathrm{cal}}$} & \multicolumn{3}{c}{ResNet-20} & \multicolumn{3}{c}{ResNet-56} \\
\cmidrule(lr){2-4} \cmidrule(lr){5-7}
 & 均值 $\bar T$ & 经验 SD & Delta SE & 均值 $\bar T$ & 经验 SD & Delta SE \\
\midrule
200 & $1.6137$ & $0.1954$ & $0.1801$ & $1.6734$ & $0.1775$ & $0.1685$ \\
300 & $1.6449$ & $0.1416$ & $0.1463$ & $1.7163$ & $0.1253$ & $0.1346$ \\
400 & $1.6776$ & $0.1133$ & $0.1266$ & $1.7403$ & $0.0784$ & $0.1150$ \\
500 & $1.6878$ & $0.0941$ & $0.1125$ & $1.7314$ & $0.0837$ & $0.1032$ \\
750 & $1.6899$ & $0.0744$ & $0.0921$ & $1.7301$ & $0.0669$ & $0.0845$ \\
1000 & $1.6760$ & $0.0820$ & $0.0797$ & $1.7302$ & $0.0622$ & $0.0729$ \\
1500 & $1.6686$ & $0.0476$ & $0.0652$ & $1.7340$ & $0.0488$ & $0.0595$ \\
2000 & $1.6666$ & $0.0418$ & $0.0566$ & $1.7272$ & $0.0331$ & $0.0516$ \\
2500 & $1.6572$ & $0.0408$ & $0.0505$ & $1.7265$ & $0.0286$ & $0.0462$ \\
3000 & $1.6670$ & $0.0407$ & $0.0461$ & $1.7289$ & $0.0329$ & $0.0421$ \\
4000 & $1.6738$ & $0.0274$ & $0.0399$ & $1.7328$ & $0.0196$ & $0.0365$ \\
\bottomrule
\end{tabular}%
}
\end{table}

\clearpage
\phantomsection
\addcontentsline{toc}{subsection}{附录 D 核心理论的 Lean 4 形式化证明源码}
\subsection*{附录 D 核心理论的 Lean 4 形式化证明源码}

本附录给出论文第~\ref{sec:theory} 节核心推导公式与定理的 Lean 4（\texttt{v4.35.0} + Mathlib \texttt{v4.35.0}）形式化证明源码（对应仓库文件 \texttt{lean\_proof/TemperatureScalingProof.lean}），已在 Lean 4 环境下通过 \texttt{lake env lean} 完整编译验证，无任何 \texttt{sorry} 或未证假设。

{\scriptsize
\begin{verbatim}
import Mathlib.Analysis.SpecialFunctions.Log.Deriv
import Mathlib.Analysis.Convex.Deriv
import Mathlib.Tactic.Linarith

open Finset Real

namespace TemperatureScaling

section SoftmaxAndOrder

variable {C : ℕ} [NeZero C]

/-- Partition function (normalizing constant) for logits `z` at inverse temperature `β`. -/
noncomputable def denom (z : Fin C → ℝ) (β : ℝ) : ℝ :=
  ∑ c : Fin C, Real.exp (β * z c)

/-- Softmax probability of class `c` at inverse temperature `β`. -/
noncomputable def softmax (z : Fin C → ℝ) (β : ℝ) (c : Fin C) : ℝ :=
  Real.exp (β * z c) / denom z β

lemma denom_pos (z : Fin C → ℝ) (β : ℝ) : 0 < denom z β := by
  unfold denom
  exact Finset.sum_pos (fun c _ ↦ Real.exp_pos (β * z c)) Finset.univ_nonempty

lemma denom_ne_zero (z : Fin C → ℝ) (β : ℝ) : denom z β ≠ 0 :=
  (denom_pos z β).ne'

lemma softmax_pos (z : Fin C → ℝ) (β : ℝ) (c : Fin C) : 0 < softmax z β c :=
  div_pos (Real.exp_pos _) (denom_pos z β)

lemma sum_softmax_eq_one (z : Fin C → ℝ) (β : ℝ) :
    ∑ c : Fin C, softmax z β c = 1 := by
  unfold softmax
  simp only [div_eq_mul_inv, ← Finset.sum_mul]
  exact mul_inv_cancel₀ (denom_ne_zero z β)

/-- Proposition 1 (Strict Order Preservation):
For `β > 0`, `z k < z c ↔ softmax z β k < softmax z β c`. -/
theorem order_preservation (z : Fin C → ℝ) {β : ℝ} (hβ : 0 < β) (c k : Fin C) :
    z k < z c ↔ softmax z β k < softmax z β c := by
  unfold softmax
  have hd : 0 < denom z β := denom_pos z β
  rw [div_lt_div_iff_of_pos_right hd, Real.exp_lt_exp, mul_lt_mul_iff_right₀ hβ]

/-- Proposition 1 Corollary (Argmax Invariance):
For `β > 0`, any maximizer of logits `z` is also a maximizer of `softmax z β`. -/
theorem argmax_invariance (z : Fin C → ℝ) {β : ℝ} (hβ : 0 < β) {c : Fin C}
    (hc : ∀ k : Fin C, z k ≤ z c) :
    ∀ k : Fin C, softmax z β k ≤ softmax z β c := by
  intro k
  unfold softmax
  have hd : 0 < denom z β := denom_pos z β
  rw [div_le_div_iff_of_pos_right hd, Real.exp_le_exp]
  exact mul_le_mul_of_nonneg_left (hc k) hβ.le

end SoftmaxAndOrder

section LikelihoodAndDerivatives

variable {C n : ℕ} [NeZero C]

/-- Conditional expectation `E_{p(β)}[Z] = ∑ c, p_c(β) z_c`. -/
noncomputable def condExpect (z : Fin C → ℝ) (β : ℝ) : ℝ :=
  ∑ c : Fin C, softmax z β c * z c

/-- Conditional variance `Var_{p(β)}(Z) = ∑ c, p_c(β) (z_c - E[Z])^2`. -/
noncomputable def condVar (z : Fin C → ℝ) (β : ℝ) : ℝ :=
  ∑ c : Fin C, softmax z β c * (z c - condExpect z β) ^ 2

/-- Algebraic identity: `∑ p_c z_c^2 - (∑ p_c z_c)^2 = ∑ p_c (z_c - E[Z])^2`. -/
theorem condVar_eq_second_moment_sub_sq (z : Fin C → ℝ) (β : ℝ) :
    (∑ c : Fin C, softmax z β c * z c ^ 2) - (condExpect z β) ^ 2 = condVar z β := by
  unfold condVar
  have hsum : ∑ c : Fin C, softmax z β c = 1 := sum_softmax_eq_one z β
  have hexp : ∑ c : Fin C, softmax z β c * z c = condExpect z β := rfl
  calc
    (∑ c : Fin C, softmax z β c * z c ^ 2) - (condExpect z β) ^ 2
      = (∑ c : Fin C, softmax z β c * z c ^ 2)
          - 2 * condExpect z β * (∑ c : Fin C, softmax z β c * z c)
          + (condExpect z β) ^ 2 * (∑ c : Fin C, softmax z β c) := by
            rw [hexp, hsum]; ring
    _ = ∑ c : Fin C, (softmax z β c * z c ^ 2
          - 2 * condExpect z β * (softmax z β c * z c)
          + (condExpect z β) ^ 2 * softmax z β c) := by
            simp only [Finset.mul_sum, Finset.sum_add_distrib, Finset.sum_sub_distrib]
    _ = ∑ c : Fin C, softmax z β c * (z c - condExpect z β) ^ 2 := by
            refine Finset.sum_congr rfl fun c _ ↦ ?_
            ring

omit [NeZero C] in
/-- Derivative of `denom z β` with respect to `β`. -/
lemma hasDerivAt_denom (z : Fin C → ℝ) (β : ℝ) :
    HasDerivAt (denom z) (∑ c : Fin C, Real.exp (β * z c) * z c) β := by
  unfold denom
  have h : ∀ c ∈ (Finset.univ : Finset (Fin C)),
      HasDerivAt (fun b ↦ Real.exp (b * z c)) (Real.exp (β * z c) * z c) β := by
    intro c _
    have hmul : HasDerivAt (fun b ↦ b * z c) (1 * z c) β :=
      (hasDerivAt_id β).mul_const (z c)
    simpa [one_mul] using hmul.exp
  exact HasDerivAt.fun_sum h

/-- Derivative of `log (denom z β)` equals `condExpect z β`. -/
lemma hasDerivAt_log_denom (z : Fin C → ℝ) (β : ℝ) :
    HasDerivAt (fun b ↦ Real.log (denom z b)) (condExpect z β) β := by
  have hd : denom z β ≠ 0 := denom_ne_zero z β
  have hlog := (hasDerivAt_denom z β).log hd
  convert hlog using 1
  unfold condExpect softmax
  simp only [div_eq_mul_inv, Finset.sum_mul]
  refine Finset.sum_congr rfl fun c _ ↦ ?_
  ring

/-- Single-sample log-likelihood `l_i(β) = β z_{y} - log(∑_c exp(β z_c))`. -/
noncomputable def sampleLogLik (z : Fin C → ℝ) (y : Fin C) (β : ℝ) : ℝ :=
  β * z y - Real.log (denom z β)

/-- Total log-likelihood `l(β) = ∑_i l_i(β)`. -/
noncomputable def totalLogLik (z : Fin n → Fin C → ℝ) (y : Fin n → Fin C) (β : ℝ) : ℝ :=
  ∑ i : Fin n, sampleLogLik (z i) (y i) β

/-- Score function `U(β) = ∑_i (z_{i, y_i} - E_{p_i(β)}[Z_i])`. -/
noncomputable def scoreFun (z : Fin n → Fin C → ℝ) (y : Fin n → Fin C) (β : ℝ) : ℝ :=
  ∑ i : Fin n, (z i (y i) - condExpect (z i) β)

/-- Score Equation Derivation: `d/dβ l(β) = U(β)`. -/
theorem hasDerivAt_totalLogLik (z : Fin n → Fin C → ℝ) (y : Fin n → Fin C) (β : ℝ) :
    HasDerivAt (totalLogLik z y) (scoreFun z y β) β := by
  unfold totalLogLik scoreFun
  have h : ∀ i ∈ (Finset.univ : Finset (Fin n)),
      HasDerivAt (sampleLogLik (z i) (y i)) (z i (y i) - condExpect (z i) β) β := by
    intro i _
    have h1 : HasDerivAt (fun b ↦ b * z i (y i)) (z i (y i)) β := by
      simpa [one_mul] using (hasDerivAt_id β).mul_const (z i (y i))
    exact h1.sub (hasDerivAt_log_denom (z i) β)
  exact HasDerivAt.fun_sum h

/-- Derivative of Softmax probability `p_c(β)` equals `p_c(β) * (z_c - E[Z])`. -/
theorem hasDerivAt_softmax (z : Fin C → ℝ) (β : ℝ) (c : Fin C) :
    HasDerivAt (fun b ↦ softmax z b c) (softmax z β c * (z c - condExpect z β)) β := by
  unfold softmax
  have hnum : HasDerivAt (fun b ↦ Real.exp (b * z c)) (Real.exp (β * z c) * z c) β := by
    have hmul : HasDerivAt (fun b ↦ b * z c) (1 * z c) β :=
      (hasDerivAt_id β).mul_const (z c)
    simpa [one_mul] using hmul.exp
  have hden := hasDerivAt_denom z β
  have hdiv := hnum.div hden (denom_ne_zero z β)
  convert hdiv using 1
  have hce : (∑ k : Fin C, Real.exp (β * z k) * z k) / denom z β = condExpect z β := by
    unfold condExpect softmax
    simp only [div_eq_mul_inv, Finset.sum_mul]
    refine Finset.sum_congr rfl fun k _ ↦ ?_
    ring
  rw [← hce]
  field_simp [denom_ne_zero z β]

/-- Derivative of `condExpect z β` equals `condVar z β`. -/
lemma hasDerivAt_condExpect (z : Fin C → ℝ) (β : ℝ) :
    HasDerivAt (condExpect z) (condVar z β) β := by
  unfold condExpect
  have h : ∀ c ∈ (Finset.univ : Finset (Fin C)),
      HasDerivAt (fun b ↦ softmax z b c * z c)
        (softmax z β c * (z c - condExpect z β) * z c) β := by
    intro c _
    exact (hasDerivAt_softmax z β c).mul_const (z c)
  have hsum := HasDerivAt.fun_sum h
  convert hsum using 1
  rw [← condVar_eq_second_moment_sub_sq z β]
  calc
    (∑ c : Fin C, softmax z β c * z c ^ 2) - (condExpect z β) ^ 2
      = (∑ c : Fin C, softmax z β c * z c ^ 2) -
          condExpect z β * (∑ c : Fin C, softmax z β c * z c) := by
            rw [show (∑ c : Fin C, softmax z β c * z c) = condExpect z β from rfl, sq]
    _ = ∑ c : Fin C, (softmax z β c * z c ^ 2 - condExpect z β * (softmax z β c * z c)) := by
            rw [Finset.mul_sum, Finset.sum_sub_distrib]
    _ = ∑ c : Fin C, softmax z β c * (z c - condExpect z β) * z c := by
            refine Finset.sum_congr rfl fun c _ ↦ ?_
            ring

/-- Observed Fisher Information `I_n(β) = ∑_i Var_{p_i(β)}(Z_i)`. -/
noncomputable def fisherInfo (z : Fin n → Fin C → ℝ) (β : ℝ) : ℝ :=
  ∑ i : Fin n, condVar (z i) β

/-- Second Derivative / Observed Fisher Information Formula: `d/dβ U(β) = - I_n(β)`. -/
theorem hasDerivAt_scoreFun (z : Fin n → Fin C → ℝ) (y : Fin n → Fin C) (β : ℝ) :
    HasDerivAt (scoreFun z y) (- fisherInfo z β) β := by
  unfold scoreFun fisherInfo
  have h : ∀ i ∈ (Finset.univ : Finset (Fin n)),
      HasDerivAt (fun b ↦ z i (y i) - condExpect (z i) b) (- condVar (z i) β) β := by
    intro i _
    have hc : HasDerivAt (fun _ : ℝ ↦ z i (y i)) 0 β := hasDerivAt_const β (z i (y i))
    simpa using hc.fun_sub (hasDerivAt_condExpect (z i) β)
  simpa [Finset.sum_neg_distrib] using HasDerivAt.fun_sum h

end LikelihoodAndDerivatives

section ConcavityAndUniqueness

variable {C n : ℕ} [NeZero C]

lemma condVar_nonneg (z : Fin C → ℝ) (β : ℝ) : 0 ≤ condVar z β := by
  unfold condVar
  exact Finset.sum_nonneg fun c _ ↦ mul_nonneg (softmax_pos z β c).le (sq_nonneg _)

/-- If logits `z` are non-degenerate (`∃ c k, z c ≠ z k`), then `0 < condVar z β`. -/
theorem condVar_pos_of_ne (z : Fin C → ℝ) (β : ℝ) (hne : ∃ c k : Fin C, z c ≠ z k) :
    0 < condVar z β := by
  rcases hne with ⟨c, k, hck⟩
  have hterm : ∃ m : Fin C, 0 < softmax z β m * (z m - condExpect z β) ^ 2 := by
    by_cases hc : z c = condExpect z β
    · refine ⟨k, mul_pos (softmax_pos z β k) (sq_pos_of_ne_zero ?_)⟩
      intro hk
      apply hck
      linarith
    · refine ⟨c, mul_pos (softmax_pos z β c) (sq_pos_of_ne_zero (sub_ne_zero.mpr hc))⟩
  rcases hterm with ⟨m, hm⟩
  unfold condVar
  exact Finset.sum_pos' (fun j _ ↦ mul_nonneg (softmax_pos z β j).le (sq_nonneg _))
    ⟨m, Finset.mem_univ m, hm⟩

/-- Theorem 1a (Strict Positivity of Observed Fisher Information):
If at least one sample has non-degenerate logits, then `I_n(β) > 0` for all `β ∈ ℝ`. -/
theorem fisherInfo_pos (z : Fin n → Fin C → ℝ) (β : ℝ)
    (hnd : ∃ (i : Fin n) (c k : Fin C), z i c ≠ z i k) :
    0 < fisherInfo z β := by
  rcases hnd with ⟨i, c, k, hck⟩
  unfold fisherInfo
  exact Finset.sum_pos' (fun j _ ↦ condVar_nonneg (z j) β)
    ⟨i, Finset.mem_univ i, condVar_pos_of_ne (z i) β ⟨c, k, hck⟩⟩

/-- Theorem 1b (Strict Concavity of Log-Likelihood):
Under non-degeneracy, `totalLogLik z y` is strictly concave on `(0, ∞)`. -/
theorem strictConcaveOn_totalLogLik (z : Fin n → Fin C → ℝ) (y : Fin n → Fin C)
    (hnd : ∃ (i : Fin n) (c k : Fin C), z i c ≠ z i k) :
    StrictConcaveOn ℝ (Set.Ioi 0) (totalLogLik z y) := by
  have h_univ : StrictConcaveOn ℝ Set.univ (totalLogLik z y) := by
    apply strictConcaveOn_of_deriv2_neg convex_univ
    · exact ContinuousOn.congr (f := totalLogLik z y)
        (fun b _ ↦ (hasDerivAt_totalLogLik z y b).continuousAt.continuousWithinAt)
        (fun _ _ ↦ rfl)
    · intro b _
      have h1 : deriv (totalLogLik z y) = scoreFun z y := by
        ext t
        exact (hasDerivAt_totalLogLik z y t).deriv
      rw [Function.iterate_succ_apply', Function.iterate_one, h1]
      have h2 := (hasDerivAt_scoreFun z y b).deriv
      rw [h2]
      exact neg_lt_zero.mpr (fisherInfo_pos z b hnd)
  exact h_univ.subset (Set.subset_univ _) (convex_Ioi 0)

/-- Theorem 1c (Global Uniqueness of MLE):
Any two critical points of `totalLogLik` on `(0, ∞)` (where `U(β) = 0`) are equal. -/
theorem mle_unique_of_score_zero (z : Fin n → Fin C → ℝ) (y : Fin n → Fin C)
    (hnd : ∃ (i : Fin n) (c k : Fin C), z i c ≠ z i k)
    {β₁ β₂ : ℝ} (_hβ₁ : 0 < β₁) (_hβ₂ : 0 < β₂)
    (h1 : scoreFun z y β₁ = 0) (h2 : scoreFun z y β₂ = 0) :
    β₁ = β₂ := by
  have hsm : StrictMono (fun b ↦ -scoreFun z y b) := by
    apply strictMono_of_deriv_pos
    intro b
    have hd : HasDerivAt (fun t ↦ -scoreFun z y t) (fisherInfo z b) b := by
      have hneg := (hasDerivAt_scoreFun z y b).fun_neg
      simpa using hneg
    rw [hd.deriv]
    exact fisherInfo_pos z b hnd
  rcases lt_trichotomy β₁ β₂ with hlt | heq | hgt
  · have := hsm hlt
    linarith
  · exact heq
  · have := hsm hgt
    linarith

end ConcavityAndUniqueness

section DeltaMethod

/-- Delta Method Derivative & Variance Formula:
For `g(β) = 1 / β`, `g'(β) = -1 / β^2`, and `[g'(β)]^2 / I = T^4 / I` where `T = 1 / β`. -/
theorem hasDerivAt_inv_temp {β : ℝ} (hβ : β ≠ 0) :
    HasDerivAt (fun b : ℝ ↦ 1 / b) (-1 / β ^ 2) β := by
  convert hasDerivAt_inv hβ using 1
  · ext b; simp [one_div]
  · ring

theorem delta_variance_identity {β I : ℝ} (hβ : β ≠ 0) :
    (-1 / β ^ 2) ^ 2 / I = (1 / β) ^ 4 / I := by
  field_simp [hβ]

end DeltaMethod

end TemperatureScaling
\end{verbatim}
}
